%-------------------------------------------------------------------------------
%	SECTION TITLE
%-------------------------------------------------------------------------------
\cvsection{Expériences professionnelles}


%-------------------------------------------------------------------------------
%	CONTENT
%-------------------------------------------------------------------------------
\begin{cventries}

	%---------------------------------------------------------
	\begin{cvexperience}
		\cvbigentry
		{Développeur Senior Symfony / VueJS} % Job title
		{\inlinelogo{src/assets/images/dinum-logo.png} DINUM} % Organization
		{Rennes} % Location
		{Juillet 2026 - Aujourd'hui} % Date(s)
		{
			\vspace{2mm}
			\begin{cvitems} % Description(s) of tasks/responsibilities
				\item {Direction interministérielle du numérique, rattachée au Premier ministre}
				\item {Pilote la transformation numérique de l'État et opère le \textbf{Réseau Interministériel de l'État} (RIE)}
			\end{cvitems}
		}

		\inlinetitle{Applications du RIE}
		\begin{cvitems}
			\item {Développement d'applications au service du \textbf{Réseau Interministériel de l'État} pour de la gestion d'incidents et de demandes.}
			\item {Développement backend \textbf{Symfony} et frontend \textbf{VueJS}}
			\item {Environnement \textbf{très contrôlé} et sécurisé : postes \textbf{offline}, sans accès Internet, avec très peu de permissions}
			\item {Adaptation des outils et des processus de développement à ces contraintes (dépendances, images, outillage disponibles hors ligne)}
			\item {Optimisation des images \textbf{Docker} : réduction de la taille et du temps de build, rationalisation des couches}
			\item {Optimisation de la \textbf{CI/CD Gitlab} : accélération des pipelines, mutualisation des jobs et mise en cache}
		\end{cvitems}

		\inlinetitle{Environnement technique}
		\vspace{-5mm}
		\begin{multicols}{2}
			\begin{cvitems}
				\item PHP 8
				\item \textbf{Symfony}
				\item \textbf{VueJS}
				\item \textbf{Docker}
				\item \textbf{Gitlab CI/CD}
				\item PHPUnit
			\end{cvitems}
		\end{multicols}
	\end{cvexperience}

	%---------------------------------------------------------
	\begin{cvexperience}
		\cvbigentry
		{Développeur Full-Stack .NET / Angular} % Job title
		{Estra Consulting} % Organization
		{Télétravail} % Location
		{Mai 2026 - Juillet 2026} % Date(s)
		{
			\vspace{2mm}
			\begin{cvitems} % Description(s) of tasks/responsibilities
				\item {Mission pour la \textbf{CAPEB} (Confédération de l'artisanat et des petites entreprises du bâtiment)}
				\item {Environ 60 000 entreprises adhérentes réparties sur les CAPEB départementales}
			\end{cvitems}
		}

		\inlinetitle{CAPEB Avantages}
		\begin{cvitems}
			\item {Site de redirection des artisans adhérents vers les offres des partenaires, et formulaire de prise de contact pour les non-adhérents}
			\item {Consommation de plusieurs \textbf{microservices} internes : référentiel adhérents et aiguillage, envoi de mails, désinscription}
			\item {\textbf{API .NET / C\#} en \textbf{architecture hexagonale}, base PostgreSQL}
			\item {Front \textbf{Angular} : validation du SIRET à la saisie, configuration externalisée pour les tests croisés entre environnements}
			\item {Sécurisation d'une API publique : \textbf{rate limiting}, CORS, CSRF}
			\item {Logs métier systématiques, endpoints de version et de santé, alignement des versions front et back}
			\item {Développement d'un mode \textit{dry-run} sur le microservice d'envoi de mails pour la préproduction}
			\item {Qualité de code contrôlée par \textbf{SonarQube} (Quality Gate bloquante), Conventional Commits et SemVer}
		\end{cvitems}

		\inlinetitle{BackOffice}
		\begin{cvitems}
			\item {Vérification et suivi des SIRET : historique des passages et des erreurs}
			\item {Suivi des mails envoyés, envoi unitaire pour le support, gestion des demandes de contact}
			\item {Configuration des partenaires et des URLs d'aiguillage}
			\item {\textbf{Statistiques} via Grafana}
		\end{cvitems}

		\inlinetitle{Environnement technique}
		\vspace{-5mm}
		\begin{multicols}{2}
			\begin{cvitems}
				\item \textbf{.NET / C\#}
				\item \textbf{Angular}
				\item PostgreSQL
				\item \textbf{Architecture hexagonale}
				\item Gitlab CI, SonarQube
				\item Grafana, Jira
			\end{cvitems}
		\end{multicols}
	\end{cvexperience}

	%---------------------------------------------------------
	\begin{cvexperience}
		\cvbigentry
		{Développeur Senior Sylius} % Job title
		{\inlinelogo{src/assets/images/antadis-logo.png} Antadis} % Organization
		{Télétravail} % Location
		{Avril 2026 - Mai 2026} % Date(s)
		{
			\vspace{2mm}
			\begin{cvitems} % Description(s) of tasks/responsibilities
				\item {Agence web spécialisée dans l'e-commerce}
				\item {Intervention en tant qu'expert Sylius / Symfony}
			\end{cvitems}
		}

		\inlinetitle{France Matériaux}
		\begin{cvitems}
			\item {Plateforme \textbf{e-commerce} fédérant les quelque 250 adhérents d'un réseau de négoce de matériaux, chacun avec ses propres pratiques}
			\item {\textbf{Sylius multi-canaux} : configuration adaptée aux besoins propres de chaque adhérent}
			\item {Gestion de \textbf{plusieurs administrateurs} et de leurs périmètres}
			\item {Recherche produit avec \textbf{ElasticSearch}, interface en \textbf{VueJS}}
			\item {Accompagnement de l'équipe sur les bonnes pratiques Sylius / Symfony}
		\end{cvitems}

		\inlinetitle{Environnement technique}
		\vspace{-5mm}
		\begin{multicols}{2}
			\begin{cvitems}
				\item \textbf{Sylius}
				\item \textbf{Symfony}
				\item \textbf{ElasticSearch}
				\item VueJS
			\end{cvitems}
		\end{multicols}
	\end{cvexperience}

	%---------------------------------------------------------
	\begin{cvexperience}
		\cvbigentry
		{Développeur Senior Symfony} % Job title
		{\inlinelogo{src/assets/images/carrefour-logo.png} Carrefour} % Organization
		{Massy} % Location
		{Juillet 2022 - Janvier 2026} % Date(s)
		{
			\vspace{2mm}
			\begin{cvitems} % Description(s) of tasks/responsibilities
				\item{Leader de la grande distribution alimentaire (magasin et livraison)}
				\item{Plus de 500 000 employés}
				\item {Chiffre d'affaires annuel de +80 milliards d'euros}
				\item {Gestion Backend du site internet de commande drive \& livraison (validation et paiement du panier, sélection de magasin et créneaux) ayant plusieurs millions de commandes/an}
				\item {Création d'une CLI pour interagir avec l'environnement de Carrefour}
			\end{cvitems}
		}

		\inlinetitle{Projet Frontone}
		\begin{cvitems} % Description(s) of tasks/responsibilities
			\item {\textbf{Référent TDD et DDD}. Enseignement des bonnes pratiques, pédagogie autour du TDD et mise en place au sein des équipes front et back.}
			\item {\textbf{Docker}, optimisation de l'environnement dev et ajouts de services}
			\item {Rattrapage de la dette technique du projet frontone, mise en place d'abstraction pour faciliter les bonnes pratiques.}
			\item {\textbf{Helm et K8S} sur \textbf{Microsoft Azure}, gestion des templates de configuration de déploiement}
			\item {Membre de l'équipe de test de l'\textbf{IA} chez Carrefour en tant que développeur, utilisation de \textbf{Gemini CLI}}
			\item {Implémentation de l'IA dans le projet frontone, définissant des agents et le fichier GEMINI.md}
			\item {\textbf{Symfony 5/6/7}, mise à jour de versions, aide au Tech lead transverse pour préparer les montées majeures}
			\item {Architecture en Microservices, structuration de la souveraineté de la donnée}
			\item {Division par 2 voire 3 suivants les requêtes du nombre d'appels aux microservices}
			\item {Optimisation des appels, \textbf{réduction de plusieurs centaines de millisecondes le traitement des appels API}}
			\item {\textbf{Augmentation de la vitesse d'implémentation} des API par une librairie interne, plusieurs centaines de lignes de code générées automatiquement par appel}
			\item {Normalisation et optimisation des DTO d'entrées/sorties de l'application}
			\item {Augmentation de la \textbf{validation systématique}, amélioration des erreurs de validation}
			\item {Amélioration de l'\textbf{architecture hexagonale}, mise en place de deptrac}
			\item {Implémentation d'un bridge entre le site et une application pour l'intégration dans \textbf{OpenAI}}
			\item {Animation de réunions et d'échanges techniques, présentation d'outils dans le cadre de la \textbf{veille technologique}}
		\end{cvitems}

		\inlinetitle{UI Library}
		\begin{cvitems}
			\item {\textbf{Référent DDD} dans l'équipe Front}
			\item {Mise en place de l'architecture hexagonale dans le contexte \textbf{VueJS}}
			\item {Aide à la mise en place de l'injection de dépendance}
			\item {Dialogue pour l'optimisation des performances entre Frontone et ce projet}
			\item {Création de l'image Docker pour les devs}
			\item {Enseignement des bonnes pratiques aux autres devs}
		\end{cvitems}

		\inlinetitle{Environnement technique}
		\vspace{-5mm}
		\begin{multicols}{2}
			\begin{cvitems}
				\item PHP 7.4, 8.0 et 8.1
				\item \textbf{Symfony 5, 6 et 7}
				\item \textbf{Docker}
				\item Vue3
				\item NodeJS
				\item Gemini CLI
				\item \textbf{Kubernetes, Azure, Helm}
				\item \textbf{DDD, TDD}
				\item Gitlab CI, Jenkins
				\item PHPUnit
			\end{cvitems}
		\end{multicols}
	\end{cvexperience}

	%---------------------------------------------------------
	\begin{cvexperience}
		\cvbigentry
		{Développeur Senior Symfony} % Job title
		{\inlinelogo{src/assets/images/icade-logo.png} Icade} % Organization
		{La Défense, Paris} % Location
		{Janvier 2022 - Juin 2022} % Date(s)
		{
			\vspace{2mm}
			\begin{cvitems} % Description(s) of tasks/responsibilities
				\item {Leader de la promotion immobilière et du foncier tertiaire}
				\item {Détenteur d'un patrimoine immobilier de plusieurs milliards d'euros}
			\end{cvitems}
		}

		\inlinetitle{Promotion interne}
		\begin{cvitems} % Description(s) of tasks/responsibilities
			\item {CMS avec un BackOffice fait par un prestataire non fonctionnel}
			\item {\textbf{Ibexa DXP} (basé sur Symfony), \textbf{Elasticsearch} et MinIO}
			\item {Nettoyage, défragmentation, sécurisation et optimisation de la stack \textbf{Docker} et images}
			\item {Nettoyage et optimisation du code du projet}
			\item {Corrections de bugs et implémentation de fonctionnalités sur le BackOffice}
			\item {Analyse statique et build des images sur une \textbf{CI Gitlab}.}
			\item {Déploiement sur Kubernetes via la CI/CD sur Gitlab}
		\end{cvitems}

		\inlinetitle{RCU (Référentiel Client Unique)}
		\begin{cvitems}
			\item {Agrégateur de données clients éparpillées sur plusieurs services internes}
			\item {Visualisation des données clients unifiées}
			\item {Symfony 6 et VueJS}
		\end{cvitems}

		\inlinetitle{Environnement technique}
		\vspace{-5mm}
		\begin{multicols}{2}
			\begin{cvitems}
				\item PHP 7.4, et 8.1
				\item \textbf{Symfony 4 et 5, Ibexa DXP}
				\item \textbf{Docker}
				\item Vue3
				\item ElasticSearch
				\item \textbf{Kubernetes, Azure}
				\item \textbf{Gitlab CI/CD}
				\item PHPUnit
			\end{cvitems}
		\end{multicols}
	\end{cvexperience}

	%---------------------------------------------------------
	\begin{cvexperience}
		\cvbigentry
		{Développeur Full-Stack} % Job title
		{\inlinelogo{src/assets/images/richid-logo.jpg} Rich-ID \& Chaplean} % Organization
		{Bordeaux} % Location
		{Août 2019 - Décembre 2021} % Date(s)
		{
			\vspace{2mm}
			\begin{cvitems} % Description(s) of tasks/responsibilities
				\item {Startup puis scaleup}
				\item {Projets : Opcalia, Fafih, Akto, HelloRSE, myCertif}
				\item {Test Driven Development (\textbf{TDD}) et Domain Driven Design (\textbf{DDD})}
				\item {Bundles \textbf{OpenSource}, \textbf{Intégration continue}, référent technique et veille technologique}
				\item {Développement d'une suite de tests OpenSource pour étendre les fonctionnalités de PHPUnit}
				\item {Développement d'un mapper de DTO d'entrées et de sorties OpenSource}
				\item {Legacy : Entretien d'un projet initialement sur Zend Framework}
				\item {Bases de données à forte volumétrie}
				\item {Contrainte de qualité exceptionnelle, demandant une grande attention sur la qualité, les tests et la relecture de code}
				\item {Prise en charge de fonctionnalité d'un point de vue théorique, production de documents de direction technique pour les autres devs.}
			\end{cvitems}
		}

		\inlinetitle{Environnement technique}
		\vspace{-5mm}
		\begin{multicols}{2}
			\begin{cvitems}
				\item PHP 7.1, 7.2, 7.3 et 7.4
				\item \textbf{Symfony 3, 4 et 5}
				\item Sylius
				\item ElasticSearch
				\item \textbf{Elm}, Twig, AngularJS
				\item \textbf{Docker}
				\item \textbf{Gitlab CI/CD}
				\item Déploiement Barebone servers
				\item PHPUnit, Behat
			\end{cvitems}
		\end{multicols}
	\end{cvexperience}

	%---------------------------------------------------------
\end{cventries}
